% !TeX root = ../papers/nc-p3.tex
\section{特異核の有限鉛筆を切る五次二項式の交代形式}
\label{sec:singular-quintic-pfaffians}
\begin{theorem}\label{thm:singular-quintic-pfaffians}
整な八余因子八次曲面の特異核で有限鉛筆元を取る。節点二捩れも含め、
$M=(za,a,zb,b):\mathcal O(-1)^4\to S_5\otimes\mathcal O$ とする。
四辺有限窓が存在するには、$M$ の階数3以下の軌跡 $Y$ が曲線を含むことが必要である。
$C=(a_0,b_0)$ として $k=\ell(X\cap C)\le3$、$\epsilon=1$ または0を
二つの残余二次式の結果式が $C$ 上で恒等零でないか否かに応じて取る。
十五個の最大小行列式の線型包の次元は $10+k+\epsilon\le14$ 以下である。
いずれも点窓の被約性を仮定しない。
\end{theorem}
\begin{proof}
$u\wedge v$ を $u(z)v(w)-v(z)u(w)$ と表す。
$(z-w)N$ は六行交代行列を与え、三表示で順に
\begin{align*}
 &za\wedge b+t\,zb\wedge a,\\
 &(za\wedge b+zb\wedge a)/t+za\wedge zb,\\
 &za\wedge a+2zb\wedge a-2za\wedge b
\end{align*}
となる。従って常に $MJM^{\mathsf t}$ であり、定数交代四行 $J$ のPfaffianは
順に $-t,-1/t^2,4$ である。この交代行列の四行Pfaffianは $M$ の対応する最大小行列式の
同じ非零定数倍である。点窓では $N$ が局所的に $yy^{\mathsf t}$ なので
$(z-w)N=zy\wedge y$ が分解可能、従って $X\subset Y$ はスキームとして成立する。

$Y$ が有限なら、最大行列式の期待余次元3における長さは
$c_3(\mathcal O^6-\mathcal O(-1)^4)=20$ である。すると $X=Y$ となる。
一方 $C$ は真の直線であり（有限非分岐では結果式、二捩れでは特異点の行による平面因子の排除）、
$C$ 上で両四次式が共通の導手二次式 $g$ で割れる。
$M$ は固定四空間 $gS_3$ に入るのでその最大小行列式の制限は、
二つの二次式 $a/g,b/g$ のSylvester行列式一つで生成される。
従って $Y\cap C$ は全直線か長さ4の四次因子であり、$\ell(X\cap C)\le3$ に反する。

より細かく $H^0(\mathcal I_X(4))$ の $C$ への像は、$X\cap C$ で消える全四次式である。
実際そのような式を入力全体に持上げた値は $X$ 上で $(a_0,b_0)$ に入り、
三次式の $X$ への評価が全射なので、二つの生成元の三次式倍を引いて $X$ 上の値を零にできる。
よってこの制限の階数は $5-k$ である。元の空間は十五次元で、$M$ の小行列式の
$C$ への制限は $\epsilon$ 次元だから求める上界が従う。
\end{proof}

\begin{proposition}\label{prop:singular-independent-common-cubic}
$X$ の任意の局所環で、剰余体上で $a,b$ が独立なら、正規化の共通三次因子と
前の各有限鉛筆の $(u,\lambda)$ が環全体で一意に復元できる（因子の定数単元倍を除く）。
とくに $M$ の退化の曲線は、導手直線か階数1軌跡の曲線である。
\end{proposition}
\begin{proof}
独立点の $X$ 上では $M$ の階数は正確に3である。二次元しか張らなければ
$S_1a=S_1b$ で二空間の最大公約因子から $a,b$ が比例してしまう。
単元小行列式で消去すると核は自由な直線、従って二本の一次式の関係 $l a+m b=0$ が得られる。
$l,m$ は剰余体で独立で、局所環上でも $S_1$ の基底となる。そのKoszul関係から
$a=mg,b=-lg$ と共通三次式 $g$ が一意に得られる。$g$ による乗法は剰余体上単射で分裂単射だから、
二変数核から $g(z)g(w)$ を除いた二行残余行列も階数1である。
その行列式の二因子のうち残余根が一致する因子は単元である。よって残る因子が零となり、
前の表示で第二残余根または持上げの係数が $u$ から決まる。
比例しないので節点の残余根は $0,\infty$ を避け、尖点では固定点0を避ける。
節点二捩れでも $m$ の根が両端に来れば剰余体で階数1に反する。
これで $u$ と係数 $\lambda$ はスキーム上で一意に戻る。

最後に $Y$ の導手外の一般台点を考える。$a,b$ が比例すれば貼合から点は $C$ に入るから、
一般台点では上と同じ三因子が存在する。$g$ は導手で消えない。
実際どちらかの切断がそこで零なら、$g$ の消滅は両切断の定数項を零にする。
二切断の貼合から一次残余の零点に、節点なら $v=t^2u$、尖点なら $1/v=1/u+2t$ が従う。
節点二捩れでも $a=(rz-s)g$ の貼合 $r g_3=s g_0$ と $b_4+b_0=a_0$ から
$b=\lambda(rz-s)g-sg$ を得る。従っていずれも残余行列式は零、$N$ は階数1である。
\end{proof}

\path{computations/verify_singular_quintic_pfaffians.py} は三表示について交代恒等式と
全十五成分でのPfaffianの倍率を有理標本で照合する。
